
\documentclass[acmsmall,screen,review,anonymous]{acmart}
\setcopyright{none}
\settopmatter{printacmref=false,printccs=false,printfolios=true}
\setlength{\headheight}{19pt}

\usepackage{booktabs}
\usepackage{tabularx}
\usepackage{longtable}
\usepackage{pdflscape}
\usepackage{array}
\usepackage{multirow}
\usepackage{amsmath}
\usepackage{xcolor}
% ============================================================
% Dedicated macro definitions for the no-metrics manuscript.
% This file replaces the project's general macros.tex.
% It contains only commands used by no-metrics.
% ============================================================

% Model-function notation:
% \model, \model_r, \model_s, s_{\model}, etc.
%
% The provide+renew pattern remains safe even when another
% package or errd.tex has already defined \model.
\providecommand{\model}{}
\renewcommand{\model}{f}

% Run-in headings used in the Introduction.
\providecommand{\parag}[1]{}
\renewcommand{\parag}[1]{\paragraph{#1}}

% Inline source-code token formatting.
\providecommand{\prog}[1]{}
\renewcommand{\prog}[1]{\texttt{#1}}

% Semantic-success subscript, e.g., E_{\sem}(U).
\providecommand{\sem}{}
\renewcommand{\sem}{\mathrm{sem}}

% Generic message-passing notation used in the appendix.
\providecommand{\Agg}{}
\renewcommand{\Agg}{\operatorname{Agg}}

\providecommand{\Msg}{}
\renewcommand{\Msg}{\operatorname{Msg}}

\providecommand{\Upd}{}
\renewcommand{\Upd}{\operatorname{Upd}}

\providecommand{\Readout}{}
\renewcommand{\Readout}{\operatorname{Readout}}

\providecommand{\ADR}{\mathrm{ADR}}
\providecommand{\ADRw}{\mathrm{ADR}^{+w}}
\providecommand{\ADRa}{\mathrm{ADR}^{+a}}
\providecommand{\SepCP}{\underline{s}^{\mathrm{CP}}}
\providecommand{\SepWilson}{\underline{s}^{\mathrm{W}}}
\providecommand{\SepCPM}{\underline{s}^{\mathrm{CP}}_{\mathrm{M}}}
\providecommand{\MarginZero}{\mathcal{M}_{0}}
\providecommand{\MarginShortcut}{\mathcal{M}_{\mathrm{sc}}}

\renewcommand{\thesection}{S\arabic{section}}
\renewcommand{\thesubsection}{S\arabic{section}.\arabic{subsection}}

\begin{document}

\title{Supplementary Material for Correct Predictions, Wrong Reasons: On the
Limits of Metrics for Measuring Reasoning Ability}

\author{Anonymous Author(s)}
\affiliation{%
  \institution{Anonymous Institution}
  \country{Anonymous}
}

\begin{abstract}
This supplementary material documents source separation, formulas,
secondary-analysis protocols, complete constraint tables, robustness
summaries, and architecture-level instantiations used in the main paper.
The matched and shortcut-account analyses use values printed to three
decimals in frozen source exports; the group analysis uses full-precision raw
prediction records.
The material deliberately treats conventional, matched, confidence-bounded,
and shortcut-adjusted quantities as observable behavioral criteria rather
than as direct measures of a latent reasoning mechanism.
\end{abstract}

\keywords{reasoning evaluation, metric validity, matched evaluation, witnesses, shortcut learning}

\maketitle

\paragraph{Organization.}
Section~S1 documents provenance and version separation. Section~S2 defines
the secondary analyses, including the full-precision group outputs.
Sections~S3--S5 report aggregate results, sensitivity analyses, and complete
constraint tables. Section~S6 provides architecture-level instantiations of
the main argument, and Section~S7 records reproduction requirements and
remaining evidence boundaries.

\section{Provenance, Versioning, and Claim Discipline}
\label{sec:supp-provenance}

The analysis uses two frozen constraint-evaluation archives and a separate
full-precision group analysis. We refer to the frozen sources as the
\emph{matched-separability archive} and the \emph{shortcut-margin archive}.
The matched-separability
archive is the sole source for row-level metric, witness, confidence-bound,
and pair-rematching analyses. The shortcut-margin archive is used only for
its ideal-floor versus residual-shortcut-ceiling boundary analysis and its
point-estimate versus confidence-bound examples.

The group analysis reads current raw prediction records separately for
3-SAT, 2-SAT, and graph coloring. A single domain-neutral implementation then
computes equal-group, equal-pair, nonlinear, variance, and repeated-success
summaries. These outputs are not numerically merged with rows from either
frozen archive.

This separation is necessary because some overlapping narrative rows differ
slightly across the two archives, consistent with reruns, parser changes, or
versioned exports. We do not average or selectively reconcile those values.
A submission artifact should replace the source-export identifiers in its
manifest with immutable repository commit identifiers and include the
full-precision per-instance predictions.

The supplied constraint tables contain \(19\) 3-SAT rows, \(59\) 2-SAT rows,
and \(45\) graph-coloring rows, for \(123\) model--size observations. The
group analysis contains \(19\), \(66\), and \(45\) model--size slices for the
same domains, respectively, for \(130\) slices in total.

\section{Secondary-Analysis Definitions}
\label{sec:supp-definitions}

\subsection{Witness Reference and Two-Sided Separability}

For SAT, the witness reference is \(\ADRw\): a positive response receives
credit only when the emitted assignment satisfies the formula. For graph
coloring, the analogous quantity is denoted \(\ADRa\): the proposed coloring
must be valid. In the frozen matched-separability table, both appear in the
same source column; this supplement writes the task-appropriate symbol
\(W\).

For each matched group \(G\), let \(\hat{\theta}_{+}(G)\) and
\(\hat{\theta}_{-}(G)\) be positive- and negative-arm recall. The point
two-sided separability is
\[
\hat{s}_G=\hat{\theta}_{+}(G)\hat{\theta}_{-}(G).
\]
The source archive reports the group-averaged point estimate together with a
factorized Clopper--Pearson lower bound \(L_{\mathrm{CP}}\) and Wilson lower
bound \(L_{\mathrm{W}}\). Because
\(\hat{s}_G=\mathrm{G\mbox{-}mean}^2\), the point statistic induces the same
classifier ordering as G-mean. The contribution of the source protocol is
the matched-group and confidence construction, not a new scalar ordering.

\subsection{Strict Ordering Discordance}

Within each task and problem size, consider model pairs for which both a
metric \(M\) and \(W\) differ at the reported precision. We compute
\[
D_M
=
\frac{
\sum_N\sum_{a<b}
\mathbf{1}\!\left[
(M_{a,N}-M_{b,N})(W_{a,N}-W_{b,N})<0
\right]
}{
\sum_N\sum_{a<b}
\mathbf{1}\!\left[
M_{a,N}\neq M_{b,N}\land W_{a,N}\neq W_{b,N}
\right]
}.
\]
No hidden precision is used to break ties. The statistic measures agreement
with a stronger observable reference; it does not assume that \(W\) is a
complete measure of reasoning.

\subsection{Metric-Level Gap}

For each row and metric \(M\), we record whether \(M>W\), \(M=W\), or \(M<W\)
at source precision and report the mean difference \(M-W\). The scales have
different operational meanings, so the difference is descriptive rather
than an estimate of ``overstatement'' in a common latent unit.

\subsection{Pair-Rematching Sensitivity}

The source tables include the original pairing, a maximum-rematching branch,
and a minimum-rematching branch. We compute
\[
\Delta_{\mathrm{match}}
=
\ADR_{\max}-\ADR_{\min}.
\]
This quantity changes only the matching rule applied to a fixed set of
predictions.

\subsection{Confidence and Boundary Sensitivity}

We summarize
\[
\Delta_{\mathrm{CI}}
=
|L_{\mathrm{CP}}-L_{\mathrm{W}}|
\]
for every constraint row. Threshold-parametric prefix boundaries are read as
\[
N^*(\tau)=
\max\left\{
N:
L_{\mathrm{CP}}(n)\ge\tau
\text{ for every tested }n\le N
\right\},
\]
subject to the source archive's prefix convention.

The shortcut-margin archive additionally defines an ideal floor
\(C_0=0.25\) and a measured residual shortcut ceiling. Its pre-registered
residual feature set is
\[
\mathcal{F}_{\mathrm{CB}}
=
\{
\texttt{mean\_clause\_len},
\texttt{min\_clause\_len},
\texttt{max\_clause\_len},
\texttt{pos\_ratio},
\texttt{var\_coverage}
\}.
\]
For each group, the archive takes the maximum two-sample KS distance over
these one-dimensional features, subtracts the \(0.95\) quantile of a
within-group label-permutation reference computed from \(200\) permutations,
and truncates at zero. Calling the resulting estimate
\(\hat d_{\mathrm{short}}\), the ceiling is
\[
\widehat C_{\mathrm{sc}}
=
\left(\frac{1+\hat d_{\mathrm{short}}}{2}\right)^2
\ge C_0.
\]
The archive's reported margin uses its own pooled factorized lower bound
\(\SepCPM\):
\[
\MarginZero=\SepCPM-C_0,
\qquad
\MarginShortcut=\SepCPM-\widehat C_{\mathrm{sc}}.
\]
We use these reported boundaries only as a design-sensitivity analysis and do
not numerically merge \(\SepCPM\) with the matched-separability archive's
group-averaged lower bound. These margins do not identify the internal
mechanism.

\subsection{Group Aggregation and Repeated-Success Outputs}
\label{sec:supp-group-definitions}

Let \(\mathcal{G}_{c}\) contain the groups with at least one complete matched
pair. For each group \(g\), let \(n_g\) be the number of complete pairs,
\(S_g\) the number whose two labels are both correct, and
\(a_g=S_g/n_g\). The equal-group and equal-pair summaries are
\[
\mathrm{GADR}
=\frac{1}{|\mathcal{G}_{c}|}\sum_{g\in\mathcal{G}_{c}}a_g,
\qquad
\mathrm{ADR}
=\frac{\sum_g S_g}{\sum_g n_g}.
\]
For \(v_g=a_g(1-a_g)\) and \(\lambda\ge0\), the analysis also exports
\[
\mathrm{SGADR}_{\lambda}
=\frac{1}{|\mathcal{G}_{c}|}
\sum_{g\in\mathcal{G}_{c}}(a_g-\lambda v_g),
\qquad
\mathrm{SGADR}_{\star,\lambda}
=\mathrm{ADR}-\lambda\,\mathrm{ADR}(1-\mathrm{ADR}).
\]
At \(\lambda=1\), these reduce to the macro-average of \(a_g^2\) and
\(\mathrm{ADR}^2\), respectively. We additionally export between-group
variance and
\[
A_{\ge k}
=\frac{1}{|\mathcal{G}_{c}|}
\sum_{g\in\mathcal{G}_{c}}\mathbf{1}[S_g\ge k],
\qquad k=1,\ldots,10.
\]

The output bundle keeps the three domains separate and includes overall,
model--size, per-group, lambda-sweep, completion, and applicability records.
The primary files contain \(19\), \(66\), and \(45\) model--size rows for
3-SAT, 2-SAT, and graph coloring. The 3-SAT and graph-coloring recomputations
match their archived summaries. For 2-SAT, the current predictions differ
from the older export in \(57\) cells across four models. The domain-neutral
implementation and the repository's original 2-SAT implementation agree on
the current inputs, so the supplement retains the older differences as a
versioning audit and reports the current-input values.

\section{Aggregate Secondary Results}
\label{sec:supp-aggregate}


\begin{table*}[t]
\centering
\caption{Metric level relative to the positive-side witness reference. Counts are computed row by row at source precision.}
\label{tab:supp-levels}
\renewcommand{\arraystretch}{1.06}
\begin{tabular}{llrrrrr}
\toprule
Domain & Metric & Rows & $M>W$ & $M=W$ & $M<W$ & Mean $(M-W)$\\
\midrule
3-SAT & Acc & 19 & 18 & 1 & 0 & 0.322 \\
3-SAT & ADR & 19 & 15 & 4 & 0 & 0.077 \\
3-SAT & $\hat{s}$ & 19 & 15 & 2 & 2 & 0.079 \\
3-SAT & $L_{\mathrm{CP}}$ & 19 & 10 & 0 & 9 & -0.010 \\
3-SAT & G-mean & 19 & 18 & 1 & 0 & 0.234 \\
3-SAT & BA & 19 & 18 & 1 & 0 & 0.323 \\
3-SAT & MCC & 19 & 14 & 2 & 3 & 0.087 \\
\midrule
2-SAT & Acc & 59 & 55 & 4 & 0 & 0.319 \\
2-SAT & ADR & 59 & 47 & 12 & 0 & 0.057 \\
2-SAT & $\hat{s}$ & 59 & 47 & 12 & 0 & 0.057 \\
2-SAT & $L_{\mathrm{CP}}$ & 59 & 32 & 1 & 26 & 0.015 \\
2-SAT & G-mean & 59 & 54 & 5 & 0 & 0.205 \\
2-SAT & BA & 59 & 55 & 4 & 0 & 0.322 \\
2-SAT & MCC & 59 & 43 & 9 & 7 & 0.069 \\
\midrule
Graph coloring & Acc & 45 & 45 & 0 & 0 & 0.490 \\
Graph coloring & ADR & 45 & 40 & 5 & 0 & 0.155 \\
Graph coloring & $\hat{s}$ & 45 & 41 & 2 & 2 & 0.159 \\
Graph coloring & $L_{\mathrm{CP}}$ & 45 & 39 & 2 & 4 & 0.114 \\
Graph coloring & G-mean & 45 & 43 & 2 & 0 & 0.359 \\
Graph coloring & BA & 45 & 45 & 0 & 0 & 0.488 \\
Graph coloring & MCC & 43 & 35 & 2 & 6 & 0.134 \\
\bottomrule
\end{tabular}
\end{table*}

\begin{table*}[t]
\centering
\caption{Strict pairwise model-ordering discordance with the witness reference. Each cell is discordant/strict-comparable pairs; source-precision ties are excluded.}
\label{tab:supp-rank-discordance}
\renewcommand{\arraystretch}{1.08}
\resizebox{\textwidth}{!}{%
\begin{tabular}{lrrrrrrrrr}
\toprule
Domain & Acc & F1$_+$ & F1$_-$ & ADR & $\hat{s}$ & $L_{\mathrm{CP}}$ & G-mean & BA & MCC \\
\midrule
3-SAT & 3/16 & 4/16 & 2/16 & 4/16 & 4/16 & 3/16 & 4/16 & 3/16 & 3/16 \\
2-SAT & 10/115 & 25/116 & 10/111 & 11/116 & 12/116 & 9/116 & 13/116 & 12/116 & 14/116 \\
Graph coloring & 25/62 & 34/62 & 21/63 & 24/62 & 26/63 & 27/63 & 30/63 & 24/62 & 25/62 \\
\bottomrule
\end{tabular}%
}
\end{table*}

\begin{table}[t]
\centering
\caption{Sensitivity to Clopper--Pearson versus Wilson lower bounds.}
\label{tab:supp-interval}
\renewcommand{\arraystretch}{1.08}
\begin{tabular}{lrrrr}
\toprule
Domain & Mean abs.\ diff. & Median & 95th pct. & Maximum\\
\midrule
3-SAT & 0.0043 & 0.0050 & 0.0073 & 0.010 \\
2-SAT & 0.0050 & 0.0060 & 0.0100 & 0.011 \\
Graph coloring & 0.0040 & 0.0030 & 0.0100 & 0.011 \\
\bottomrule
\end{tabular}
\end{table}

\begin{table}[t]
\centering
\caption{Pair-rematching sensitivity, $\Delta_{\mathrm{match}}=\mathrm{ADR}_{\max}-\mathrm{ADR}_{\min}$.}
\label{tab:supp-pair-summary}
\renewcommand{\arraystretch}{1.08}
\begin{tabular}{lrrrrr}
\toprule
Domain & Rows & Mean & Median & Maximum & Nonzero rows\\
\midrule
3-SAT & 19 & 0.026 & 0.021 & 0.080 & 78.9\% \\
2-SAT & 59 & 0.023 & 0.000 & 0.182 & 42.4\% \\
Graph coloring & 45 & 0.072 & 0.019 & 0.364 & 73.3\% \\
\bottomrule
\end{tabular}
\end{table}

\begin{table*}[t]
\centering
\caption{Largest pair-rematching spreads in each domain.}
\label{tab:supp-pair-top}
\renewcommand{\arraystretch}{1.06}
\begin{tabular}{llrrrrr}
\toprule
Domain & Model & $N$ & ADR$_{\mathrm{pair}}$ & ADR$_{\max}$ & ADR$_{\min}$ & $\Delta_{\mathrm{match}}$\\
\midrule
3-SAT & DeepSeek-Reasoner & 50 & 0.090 & 0.110 & 0.030 & 0.080 \\
3-SAT & Opus 4.7 (med. think) & 50 & 0.200 & 0.220 & 0.140 & 0.080 \\
3-SAT & Opus 4.7 (med. think) & 60 & 0.156 & 0.169 & 0.117 & 0.052 \\
3-SAT & DeepSeek-Reasoner & 60 & 0.086 & 0.086 & 0.038 & 0.048 \\
3-SAT & DeepSeek-Reasoner & 25 & 0.289 & 0.299 & 0.258 & 0.041 \\
\midrule
2-SAT & gpt-oss-20b (reason) & 25 & 0.227 & 0.364 & 0.182 & 0.182 \\
2-SAT & Qwen3-14B (think) & 50 & 0.154 & 0.231 & 0.063 & 0.168 \\
2-SAT & Qwen3-14B (think) & 40 & 0.203 & 0.238 & 0.077 & 0.161 \\
2-SAT & Qwen3-32B (think) & 50 & 0.143 & 0.185 & 0.042 & 0.143 \\
2-SAT & gpt-oss-20b (reason) & 10 & 0.727 & 0.818 & 0.682 & 0.136 \\
\midrule
Graph coloring & Llama-Nemo-49B (r-on) & 60 & 0.000 & 0.364 & 0.000 & 0.364 \\
Graph coloring & Llama-Nemo-49B (r-on) & 50 & 0.282 & 0.473 & 0.191 & 0.282 \\
Graph coloring & QwQ-32B (think) & 40 & 0.309 & 0.436 & 0.155 & 0.281 \\
Graph coloring & Llama-Nemo-49B (r-on) & 40 & 0.136 & 0.273 & 0.000 & 0.273 \\
Graph coloring & Qwen3-32B (think) & 20 & 0.278 & 0.389 & 0.141 & 0.248 \\
\bottomrule
\end{tabular}
\end{table*}

\subsection{Selected Strict Ordering Reversals}

\begin{table*}[t]
\centering
\caption{Representative model-order reversals relative to the witness
reference. Arrows indicate the ordering induced by the displayed metric and
by \(W\).}
\label{tab:supp-reversal-examples}
\renewcommand{\arraystretch}{1.08}
\begin{tabularx}{\textwidth}{p{0.16\textwidth}p{0.30\textwidth}p{0.25\textwidth}X}
\toprule
Domain/size & Models & Output-metric ordering & Witness ordering\\
\midrule
3-SAT, \(N=25\)
& Opus 4.7 vs.\ GPT-5
& Acc \(0.935>0.864\); ADR \(0.759>0.736\);
\(\hat{s}\;0.783>0.732\)
& \(W_{\mathrm{Opus}}=0.655 <
W_{\mathrm{GPT}}=0.673\)\\
Graph coloring, \(N=10\)
& Llama-Nemo-49B vs.\ GPT-OSS-20B
& Acc \(0.909>0.886\);
\(L_{\mathrm{CP}}\;0.751>0.710\)
& \(W_{\mathrm{Llama}}=0.227 <
W_{\mathrm{GPT\mbox{-}OSS}}=0.818\)\\
Graph coloring, \(N=8\)
& Llama-Nemo-49B vs.\ Qwen3-32B
& Acc \(0.932>0.824\); ADR \(0.864>0.647\)
& \(W_{\mathrm{Llama}}=0.136 <
W_{\mathrm{Qwen}}=0.482\)\\
2-SAT, \(N=30\)
& Qwen3-32B vs.\ Llama-Nemotron-Super-49B (reason)
& ADR \(0.157<0.250\); \(\hat{s}\;0.156<0.182\)
& \(W_{\mathrm{Qwen}}=0.083 >
W_{\mathrm{Llama}}=0.000\)\\
\bottomrule
\end{tabularx}
\end{table*}

\section{Threshold and Shortcut-Account Sensitivity}
\label{sec:supp-boundaries}

\subsection{Threshold-Parametric Matched-Separability Boundaries}

\begin{table}[t]
\centering
\caption{3-SAT prefix boundaries \(N^*(\tau)\) from the
matched-separability archive.}
\label{tab:supp-threshold-3sat}
\renewcommand{\arraystretch}{1.08}
\begin{tabular}{lrrrr}
\toprule
Model & \(N^*(0.50)\) & \(N^*(0.25)\) & \(N^*(0.10)\) & \(N^*(0.05)\)\\
\midrule
GPT-5 & 25 & 30 & 40 & 60\\
Opus 4.7 medium-thinking & 25 & 25 & 50 & 60\\
DeepSeek-Reasoner & 10 & 10 & 25 & 25\\
\bottomrule
\end{tabular}
\end{table}

\begin{table}[t]
\centering
\caption{2-SAT prefix boundaries \(N^*(\tau)\) from the
matched-separability archive.}
\label{tab:supp-threshold-2sat}
\renewcommand{\arraystretch}{1.08}
\begin{tabular}{lrrrr}
\toprule
Model & \(N^*(0.50)\) & \(N^*(0.25)\) & \(N^*(0.10)\) & \(N^*(0.05)\)\\
\midrule
Qwen3-14B thinking & 15 & 20 & 50 & 50\\
Qwen3-32B thinking & 10 & 15 & 20 & 30\\
QwQ-32B thinking & 10 & 15 & 20 & 25\\
\bottomrule
\end{tabular}
\end{table}

\subsection{Ideal-Floor Versus Residual-Shortcut Boundaries}

\begin{table*}[t]
\centering
\caption{Positive-margin boundaries from the shortcut-margin archive.
\(\MarginZero\) uses the ideal \(0.25\) floor; \(\MarginShortcut\) charges the measured
residual structural shortcut ceiling. A dash means no positive boundary on
the tested grid.}
\label{tab:supp-rmc-boundaries}
\renewcommand{\arraystretch}{1.08}
\begin{tabular}{llrrr}
\toprule
Task & Model & \(\MarginZero\) boundary & \(\MarginShortcut\) boundary & Shift\\
\midrule
3-SAT & GPT-5 & 30 & 25 & \(-5\)\\
3-SAT & Opus 4.7 medium-thinking & 25 & 25 & \(0\)\\
3-SAT & DeepSeek-Reasoner & 10 & 10 & \(0\)\\
\midrule
2-SAT & Qwen3-32B thinking & 15 & 10 & \(-5\)\\
2-SAT & gpt-oss-20b reasoning & 20 & 15 & \(-5\)\\
2-SAT & Nemotron-H-47B (r-on) & -- & -- & --\\
2-SAT & LN-Super-49B (r-on) & 10 & 8 & \(-2\)\\
\midrule
Graph coloring & Qwen3-32B thinking & 10 & 8 & \(-2\)\\
Graph coloring & gpt-oss-20b reasoning & 10 & 10 & \(0\)\\
Graph coloring & Nemotron-H-47B (r-on) & 15 & 10 & \(-5\)\\
Graph coloring & LN-Super-49B (r-on) & 10 & 10 & \(0\)\\
\bottomrule
\end{tabular}
\end{table*}

\subsection{Point Estimate Versus Confidence Requirement}

\begin{table}[t]
\centering
\caption{Examples in which an ideal-floor point margin is positive but its
Clopper--Pearson lower-bound margin is negative.}
\label{tab:supp-point-cp}
\renewcommand{\arraystretch}{1.08}
\begin{tabular}{llrr}
\toprule
Task & Model/size & Point \(\MarginZero\) & CP \(\MarginZero\)\\
\midrule
3-SAT & GPT-5, \(N=40\) & \(+0.053\) & \(-0.040\)\\
2-SAT & gpt-oss-20b, \(N=25\) & \(+0.156\) & \(-0.045\)\\
\bottomrule
\end{tabular}
\end{table}

\section{Complete Constraint Tables}
\label{sec:supp-complete}


\begin{landscape}
\scriptsize
\setlength{\tabcolsep}{2.7pt}
\renewcommand{\arraystretch}{1.05}
\begin{longtable}{p{4.1cm}r*{13}{r}}
\caption{Complete 3-SAT model--size table from the matched-separability archive. $W$ is the positive-side witness-aware rate: $\mathrm{ADR}^{+w}$ for SAT and $\mathrm{ADR}^{+a}$ for graph coloring. Values are reproduced at the three-decimal precision of the source export.}\label{tab:supp-full-3sat}\\
\toprule
Model & $N$ & Acc & F1$_+$ & F1$_-$ & ADR & $W$ & $\hat{s}$ & $L_{\mathrm{CP}}$ & $L_{\mathrm{W}}$ & BA & MCC & ADR$_{\mathrm{pair}}$ & ADR$_{\max}$ & ADR$_{\min}$ \\
\midrule
\endfirsthead
\multicolumn{15}{c}{\tablename\ \thetable\ continued}\\
\toprule
Model & $N$ & Acc & F1$_+$ & F1$_-$ & ADR & $W$ & $\hat{s}$ & $L_{\mathrm{CP}}$ & $L_{\mathrm{W}}$ & BA & MCC & ADR$_{\mathrm{pair}}$ & ADR$_{\max}$ & ADR$_{\min}$ \\
\midrule
\endhead
\midrule \multicolumn{15}{r}{Continued on next page}\\
\endfoot
\bottomrule
\endlastfoot
GPT-5 & 5 & 1.000 & 1.000 & 1.000 & 1.000 & 1.000 & 1.000 & 0.929 & 0.927 & 1.000 & 1.000 & 1.000 & 1.000 & 1.000 \\
GPT-5 & 10 & 0.990 & 0.991 & 0.989 & 0.989 & 0.989 & 0.981 & 0.897 & 0.897 & 0.991 & 0.980 & 0.989 & 0.989 & 0.989 \\
GPT-5 & 25 & 0.864 & 0.847 & 0.877 & 0.736 & 0.673 & 0.732 & 0.612 & 0.615 & 0.864 & 0.745 & 0.736 & 0.755 & 0.727 \\
GPT-5 & 30 & 0.735 & 0.651 & 0.787 & 0.486 & 0.440 & 0.480 & 0.367 & 0.372 & 0.734 & 0.533 & 0.486 & 0.495 & 0.468 \\
GPT-5 & 40 & 0.645 & 0.466 & 0.735 & 0.309 & 0.164 & 0.303 & 0.210 & 0.216 & 0.645 & 0.393 & 0.309 & 0.309 & 0.291 \\
GPT-5 & 50 & 0.573 & 0.288 & 0.695 & 0.173 & 0.064 & 0.168 & 0.099 & 0.105 & 0.573 & 0.242 & 0.173 & 0.173 & 0.145 \\
GPT-5 & 60 & 0.555 & 0.222 & 0.688 & 0.127 & 0.000 & 0.125 & 0.067 & 0.072 & 0.555 & 0.210 & 0.127 & 0.127 & 0.118 \\
Opus 4.7 (med. think) & 5 & 0.945 & 0.947 & 0.942 & 0.897 & 0.813 & 0.893 & 0.782 & 0.784 & 0.945 & 0.892 & 0.897 & 0.907 & 0.888 \\
Opus 4.7 (med. think) & 10 & 0.959 & 0.960 & 0.958 & 0.917 & 0.861 & 0.919 & 0.816 & 0.817 & 0.959 & 0.918 & 0.917 & 0.926 & 0.917 \\
Opus 4.7 (med. think) & 25 & 0.935 & 0.956 & 0.875 & 0.759 & 0.655 & 0.783 & 0.592 & 0.602 & 0.894 & 0.838 & 0.759 & 0.759 & 0.759 \\
Opus 4.7 (med. think) & 50 & 0.527 & 0.620 & 0.372 & 0.200 & 0.060 & 0.230 & 0.124 & 0.131 & 0.555 & 0.134 & 0.200 & 0.220 & 0.140 \\
Opus 4.7 (med. think) & 60 & 0.530 & 0.656 & 0.256 & 0.156 & 0.000 & 0.152 & 0.078 & 0.084 & 0.549 & 0.157 & 0.156 & 0.169 & 0.117 \\
DeepSeek-Reasoner & 5 & 0.986 & 0.986 & 0.987 & 0.973 & 0.973 & 0.973 & 0.892 & 0.892 & 0.986 & 0.973 & 0.973 & 0.973 & 0.973 \\
DeepSeek-Reasoner & 10 & 0.876 & 0.865 & 0.886 & 0.780 & 0.780 & 0.761 & 0.641 & 0.644 & 0.879 & 0.774 & 0.780 & 0.780 & 0.770 \\
DeepSeek-Reasoner & 25 & 0.626 & 0.454 & 0.716 & 0.289 & 0.144 & 0.290 & 0.197 & 0.203 & 0.633 & 0.349 & 0.289 & 0.299 & 0.258 \\
DeepSeek-Reasoner & 30 & 0.515 & 0.154 & 0.660 & 0.073 & 0.010 & 0.082 & 0.036 & 0.041 & 0.519 & 0.075 & 0.073 & 0.073 & 0.052 \\
DeepSeek-Reasoner & 40 & 0.524 & 0.183 & 0.664 & 0.093 & 0.031 & 0.104 & 0.048 & 0.053 & 0.520 & 0.072 & 0.093 & 0.103 & 0.082 \\
DeepSeek-Reasoner & 50 & 0.486 & 0.194 & 0.622 & 0.090 & 0.000 & 0.102 & 0.052 & 0.057 & 0.496 & -0.011 & 0.090 & 0.110 & 0.030 \\
DeepSeek-Reasoner & 60 & 0.509 & 0.160 & 0.653 & 0.086 & 0.010 & 0.086 & 0.039 & 0.044 & 0.513 & 0.049 & 0.086 & 0.086 & 0.038 \\
\end{longtable}
\end{landscape}

\begin{landscape}
\scriptsize
\setlength{\tabcolsep}{2.7pt}
\renewcommand{\arraystretch}{1.05}
\begin{longtable}{p{4.1cm}r*{13}{r}}
\caption{Complete 2-SAT model--size table from the matched-separability archive. $W$ is the positive-side witness-aware rate: $\mathrm{ADR}^{+w}$ for SAT and $\mathrm{ADR}^{+a}$ for graph coloring. Values are reproduced at the three-decimal precision of the source export.}\label{tab:supp-full-2sat}\\
\toprule
Model & $N$ & Acc & F1$_+$ & F1$_-$ & ADR & $W$ & $\hat{s}$ & $L_{\mathrm{CP}}$ & $L_{\mathrm{W}}$ & BA & MCC & ADR$_{\mathrm{pair}}$ & ADR$_{\max}$ & ADR$_{\min}$ \\
\midrule
\endfirsthead
\multicolumn{15}{c}{\tablename\ \thetable\ continued}\\
\toprule
Model & $N$ & Acc & F1$_+$ & F1$_-$ & ADR & $W$ & $\hat{s}$ & $L_{\mathrm{CP}}$ & $L_{\mathrm{W}}$ & BA & MCC & ADR$_{\mathrm{pair}}$ & ADR$_{\max}$ & ADR$_{\min}$ \\
\midrule
\endhead
\midrule \multicolumn{15}{r}{Continued on next page}\\
\endfoot
\bottomrule
\endlastfoot
gpt-oss-20b (reason) & 3 & 1.000 & 1.000 & 1.000 & 1.000 & 1.000 & 1.000 & 0.976 & 0.969 & 1.000 & 1.000 & 1.000 & 1.000 & 1.000 \\
gpt-oss-20b (reason) & 5 & 1.000 & 1.000 & 1.000 & 1.000 & 1.000 & 1.000 & 0.976 & 0.969 & 1.000 & 1.000 & 1.000 & 1.000 & 1.000 \\
gpt-oss-20b (reason) & 8 & 0.955 & 0.952 & 0.957 & 0.909 & 0.864 & 0.909 & 0.863 & 0.855 & 0.955 & 0.913 & 0.909 & 0.909 & 0.909 \\
gpt-oss-20b (reason) & 10 & 0.841 & 0.844 & 0.837 & 0.727 & 0.682 & 0.707 & 0.631 & 0.620 & 0.841 & 0.683 & 0.727 & 0.818 & 0.682 \\
gpt-oss-20b (reason) & 15 & 0.773 & 0.737 & 0.800 & 0.545 & 0.500 & 0.562 & 0.485 & 0.475 & 0.773 & 0.567 & 0.545 & 0.636 & 0.545 \\
gpt-oss-20b (reason) & 20 & 0.773 & 0.706 & 0.815 & 0.545 & 0.409 & 0.545 & 0.467 & 0.457 & 0.773 & 0.612 & 0.545 & 0.545 & 0.545 \\
gpt-oss-20b (reason) & 25 & 0.568 & 0.457 & 0.642 & 0.227 & 0.182 & 0.281 & 0.216 & 0.210 & 0.568 & 0.149 & 0.227 & 0.364 & 0.182 \\
gpt-oss-20b (reason) & 30 & 0.568 & 0.296 & 0.689 & 0.136 & 0.045 & 0.169 & 0.118 & 0.116 & 0.568 & 0.215 & 0.136 & 0.182 & 0.136 \\
gpt-oss-20b (reason) & 40 & 0.545 & 0.286 & 0.667 & 0.182 & 0.000 & 0.165 & 0.112 & 0.110 & 0.545 & 0.132 & 0.182 & 0.182 & 0.091 \\
gpt-oss-20b (reason) & 50 & 0.500 & 0.083 & 0.656 & 0.045 & 0.000 & 0.045 & 0.026 & 0.026 & 0.500 & 0.000 & 0.045 & 0.045 & 0.045 \\
Qwen3-14B (think) & 3 & 0.990 & 0.990 & 0.990 & 0.981 & 0.981 & 0.981 & 0.950 & 0.943 & 0.990 & 0.981 & 0.981 & 0.981 & 0.981 \\
Qwen3-14B (think) & 5 & 0.974 & 0.973 & 0.975 & 0.948 & 0.942 & 0.948 & 0.909 & 0.901 & 0.974 & 0.949 & 0.948 & 0.948 & 0.948 \\
Qwen3-14B (think) & 8 & 0.974 & 0.973 & 0.975 & 0.948 & 0.942 & 0.948 & 0.907 & 0.899 & 0.974 & 0.949 & 0.948 & 0.948 & 0.948 \\
Qwen3-14B (think) & 10 & 0.942 & 0.938 & 0.945 & 0.883 & 0.877 & 0.883 & 0.833 & 0.824 & 0.942 & 0.889 & 0.883 & 0.883 & 0.883 \\
Qwen3-14B (think) & 15 & 0.844 & 0.815 & 0.865 & 0.688 & 0.656 & 0.688 & 0.620 & 0.611 & 0.844 & 0.724 & 0.688 & 0.688 & 0.688 \\
Qwen3-14B (think) & 20 & 0.669 & 0.519 & 0.748 & 0.338 & 0.286 & 0.348 & 0.282 & 0.276 & 0.669 & 0.432 & 0.338 & 0.357 & 0.338 \\
Qwen3-14B (think) & 25 & 0.604 & 0.371 & 0.711 & 0.214 & 0.156 & 0.230 & 0.179 & 0.175 & 0.604 & 0.309 & 0.214 & 0.234 & 0.208 \\
Qwen3-14B (think) & 30 & 0.516 & 0.280 & 0.636 & 0.162 & 0.058 & 0.153 & 0.105 & 0.103 & 0.516 & 0.043 & 0.162 & 0.188 & 0.078 \\
Qwen3-14B (think) & 40 & 0.490 & 0.338 & 0.585 & 0.203 & 0.000 & 0.184 & 0.136 & 0.133 & 0.495 & -0.012 & 0.203 & 0.238 & 0.077 \\
Qwen3-14B (think) & 50 & 0.479 & 0.313 & 0.580 & 0.154 & 0.000 & 0.163 & 0.118 & 0.115 & 0.479 & -0.048 & 0.154 & 0.231 & 0.063 \\
Qwen3-32B (think) & 3 & 0.979 & 0.979 & 0.980 & 0.959 & 0.959 & 0.959 & 0.922 & 0.914 & 0.979 & 0.959 & 0.959 & 0.959 & 0.959 \\
Qwen3-32B (think) & 5 & 0.955 & 0.952 & 0.957 & 0.909 & 0.909 & 0.909 & 0.864 & 0.855 & 0.955 & 0.913 & 0.909 & 0.909 & 0.909 \\
Qwen3-32B (think) & 8 & 0.926 & 0.920 & 0.931 & 0.851 & 0.843 & 0.851 & 0.798 & 0.790 & 0.926 & 0.861 & 0.851 & 0.851 & 0.851 \\
Qwen3-32B (think) & 10 & 0.905 & 0.895 & 0.913 & 0.810 & 0.777 & 0.810 & 0.753 & 0.744 & 0.905 & 0.825 & 0.810 & 0.810 & 0.810 \\
Qwen3-32B (think) & 15 & 0.748 & 0.663 & 0.799 & 0.496 & 0.430 & 0.496 & 0.422 & 0.413 & 0.748 & 0.574 & 0.496 & 0.496 & 0.496 \\
Qwen3-32B (think) & 20 & 0.624 & 0.397 & 0.727 & 0.248 & 0.207 & 0.248 & 0.189 & 0.185 & 0.624 & 0.376 & 0.248 & 0.248 & 0.248 \\
Qwen3-32B (think) & 25 & 0.568 & 0.246 & 0.698 & 0.142 & 0.092 & 0.140 & 0.098 & 0.096 & 0.567 & 0.254 & 0.142 & 0.142 & 0.133 \\
Qwen3-32B (think) & 30 & 0.574 & 0.270 & 0.700 & 0.157 & 0.083 & 0.156 & 0.112 & 0.110 & 0.574 & 0.270 & 0.157 & 0.157 & 0.149 \\
Qwen3-32B (think) & 40 & 0.525 & 0.173 & 0.667 & 0.099 & 0.000 & 0.096 & 0.067 & 0.066 & 0.525 & 0.094 & 0.099 & 0.099 & 0.083 \\
Qwen3-32B (think) & 50 & 0.467 & 0.273 & 0.579 & 0.143 & 0.000 & 0.129 & 0.089 & 0.087 & 0.464 & -0.083 & 0.143 & 0.185 & 0.042 \\
QwQ-32B (think) & 3 & 1.000 & 1.000 & 1.000 & 1.000 & 1.000 & 1.000 & 0.976 & 0.969 & 1.000 & 1.000 & 1.000 & 1.000 & 1.000 \\
QwQ-32B (think) & 5 & 0.990 & 0.990 & 0.990 & 0.980 & 0.980 & 0.980 & 0.949 & 0.942 & 0.990 & 0.980 & 0.980 & 0.980 & 0.980 \\
QwQ-32B (think) & 8 & 0.980 & 0.980 & 0.980 & 0.960 & 0.960 & 0.960 & 0.924 & 0.916 & 0.980 & 0.960 & 0.960 & 0.960 & 0.960 \\
QwQ-32B (think) & 10 & 0.949 & 0.947 & 0.952 & 0.899 & 0.869 & 0.899 & 0.857 & 0.849 & 0.949 & 0.904 & 0.899 & 0.899 & 0.899 \\
QwQ-32B (think) & 15 & 0.798 & 0.747 & 0.832 & 0.596 & 0.545 & 0.596 & 0.521 & 0.511 & 0.798 & 0.652 & 0.596 & 0.596 & 0.596 \\
QwQ-32B (think) & 20 & 0.641 & 0.441 & 0.736 & 0.283 & 0.242 & 0.283 & 0.222 & 0.218 & 0.641 & 0.406 & 0.283 & 0.283 & 0.283 \\
QwQ-32B (think) & 25 & 0.576 & 0.276 & 0.700 & 0.162 & 0.111 & 0.161 & 0.122 & 0.119 & 0.576 & 0.270 & 0.162 & 0.162 & 0.152 \\
QwQ-32B (think) & 30 & 0.515 & 0.059 & 0.673 & 0.030 & 0.000 & 0.030 & 0.017 & 0.017 & 0.515 & 0.124 & 0.030 & 0.030 & 0.030 \\
QwQ-32B (think) & 40 & 0.487 & 0.093 & 0.642 & 0.033 & 0.011 & 0.039 & 0.018 & 0.019 & 0.509 & 0.043 & 0.033 & 0.044 & 0.022 \\
QwQ-32B (think) & 50 & 0.506 & 0.044 & 0.667 & 0.023 & 0.000 & 0.023 & 0.013 & 0.013 & 0.506 & 0.044 & 0.023 & 0.023 & 0.023 \\
Llama-Nemotron-Super-49B (reason) & 3 & 1.000 & 1.000 & 1.000 & 1.000 & 0.909 & 1.000 & 0.976 & 0.969 & 1.000 & 1.000 & 1.000 & 1.000 & 1.000 \\
Llama-Nemotron-Super-49B (reason) & 5 & 1.000 & 1.000 & 1.000 & 1.000 & 1.000 & 1.000 & 0.976 & 0.969 & 1.000 & 1.000 & 1.000 & 1.000 & 1.000 \\
Llama-Nemotron-Super-49B (reason) & 8 & 0.864 & 0.842 & 0.880 & 0.727 & 0.727 & 0.727 & 0.653 & 0.642 & 0.864 & 0.756 & 0.727 & 0.727 & 0.727 \\
Llama-Nemotron-Super-49B (reason) & 10 & 0.773 & 0.706 & 0.815 & 0.545 & 0.545 & 0.545 & 0.467 & 0.457 & 0.773 & 0.612 & 0.545 & 0.545 & 0.545 \\
Llama-Nemotron-Super-49B (reason) & 15 & 0.727 & 0.625 & 0.786 & 0.455 & 0.364 & 0.455 & 0.377 & 0.369 & 0.727 & 0.542 & 0.455 & 0.455 & 0.455 \\
Llama-Nemotron-Super-49B (reason) & 20 & 0.818 & 0.778 & 0.846 & 0.636 & 0.545 & 0.636 & 0.558 & 0.548 & 0.818 & 0.683 & 0.636 & 0.636 & 0.636 \\
Llama-Nemotron-Super-49B (reason) & 25 & 0.773 & 0.706 & 0.815 & 0.545 & 0.182 & 0.545 & 0.467 & 0.457 & 0.773 & 0.612 & 0.545 & 0.545 & 0.545 \\
Llama-Nemotron-Super-49B (reason) & 30 & 0.400 & 0.308 & 0.471 & 0.250 & 0.000 & 0.182 & 0.094 & 0.095 & 0.591 & 0.237 & 0.250 & 0.250 & 0.250 \\
Nemotron-H-47B-Reasoning (reason) & 3 & 0.524 & 0.657 & 0.224 & 0.077 & 0.046 & 0.080 & 0.069 & 0.068 & 0.524 & 0.076 & 0.077 & 0.087 & 0.074 \\
Nemotron-H-47B-Reasoning (reason) & 5 & 0.510 & 0.668 & 0.066 & 0.023 & 0.004 & 0.027 & 0.019 & 0.019 & 0.510 & 0.062 & 0.023 & 0.031 & 0.021 \\
Nemotron-H-47B-Reasoning (reason) & 8 & 0.499 & 0.666 & -- & 0.000 & 0.000 & 0.000 & 0.000 & 0.000 & 0.499 & -0.031 & 0.000 & 0.000 & 0.000 \\
Nemotron-H-47B-Reasoning (reason) & 10 & 0.497 & 0.664 & 0.004 & 0.002 & 0.000 & 0.002 & 0.001 & 0.001 & 0.497 & -0.042 & 0.002 & 0.002 & 0.002 \\
Nemotron-H-47B-Reasoning (reason) & 15 & 0.513 & 0.668 & 0.087 & 0.027 & 0.000 & 0.030 & 0.024 & 0.024 & 0.513 & 0.070 & 0.027 & 0.031 & 0.025 \\
Nemotron-H-47B-Reasoning (reason) & 20 & 0.510 & 0.664 & 0.093 & 0.043 & 0.000 & 0.043 & 0.037 & 0.036 & 0.510 & 0.049 & 0.043 & 0.046 & 0.041 \\
Nemotron-H-47B-Reasoning (reason) & 25 & 0.526 & 0.666 & 0.183 & 0.070 & 0.000 & 0.071 & 0.058 & 0.057 & 0.526 & 0.096 & 0.070 & 0.079 & 0.062 \\
Nemotron-H-47B-Reasoning (reason) & 30 & 0.503 & 0.665 & 0.034 & 0.015 & 0.000 & 0.017 & 0.012 & 0.012 & 0.503 & 0.024 & 0.015 & 0.017 & 0.015 \\
Nemotron-H-47B-Reasoning (reason) & 40 & 0.527 & 0.673 & 0.147 & 0.074 & 0.000 & 0.074 & 0.055 & 0.054 & 0.527 & 0.120 & 0.074 & 0.081 & 0.062 \\
Nemotron-H-47B-Reasoning (reason) & 50 & 0.516 & 0.667 & 0.120 & 0.046 & 0.000 & 0.052 & 0.042 & 0.041 & 0.516 & 0.076 & 0.046 & 0.062 & 0.041 \\
Nemotron-H-47B-Reasoning (reason) & 60 & 0.537 & 0.674 & 0.206 & 0.105 & 0.000 & 0.107 & 0.085 & 0.083 & 0.535 & 0.125 & 0.105 & 0.119 & 0.093 \\
\end{longtable}
\end{landscape}

\begin{landscape}
\scriptsize
\setlength{\tabcolsep}{2.7pt}
\renewcommand{\arraystretch}{1.05}
\begin{longtable}{p{4.1cm}r*{13}{r}}
\caption{Complete graph-coloring model--size table from the matched-separability archive. $W$ is the positive-side witness-aware rate: $\mathrm{ADR}^{+w}$ for SAT and $\mathrm{ADR}^{+a}$ for graph coloring. Values are reproduced at the three-decimal precision of the source export.}\label{tab:supp-full-graph}\\
\toprule
Model & $N$ & Acc & F1$_+$ & F1$_-$ & ADR & $W$ & $\hat{s}$ & $L_{\mathrm{CP}}$ & $L_{\mathrm{W}}$ & BA & MCC & ADR$_{\mathrm{pair}}$ & ADR$_{\max}$ & ADR$_{\min}$ \\
\midrule
\endfirsthead
\multicolumn{15}{c}{\tablename\ \thetable\ continued}\\
\toprule
Model & $N$ & Acc & F1$_+$ & F1$_-$ & ADR & $W$ & $\hat{s}$ & $L_{\mathrm{CP}}$ & $L_{\mathrm{W}}$ & BA & MCC & ADR$_{\mathrm{pair}}$ & ADR$_{\max}$ & ADR$_{\min}$ \\
\midrule
\endhead
\midrule \multicolumn{15}{r}{Continued on next page}\\
\endfoot
\bottomrule
\endlastfoot
GPT-OSS-20B (think) & 8 & 0.953 & 0.952 & 0.955 & 0.905 & 0.857 & 0.909 & 0.852 & 0.843 & 0.955 & 0.911 & 0.905 & 0.905 & 0.905 \\
GPT-OSS-20B (think) & 10 & 0.886 & 0.878 & 0.894 & 0.818 & 0.818 & 0.781 & 0.710 & 0.700 & 0.886 & 0.780 & 0.818 & 0.818 & 0.773 \\
GPT-OSS-20B (think) & 15 & 0.545 & 0.286 & 0.667 & 0.182 & 0.136 & 0.165 & 0.130 & 0.127 & 0.545 & 0.132 & 0.182 & 0.182 & 0.136 \\
GPT-OSS-20B (think) & 20 & 0.659 & 0.462 & 0.750 & 0.316 & 0.316 & 0.313 & 0.242 & 0.236 & 0.650 & 0.424 & 0.322 & 0.322 & 0.322 \\
GPT-OSS-20B (think) & 30 & 0.528 & -- & 0.691 & 0.000 & 0.000 & 0.000 & 0.000 & 0.000 & 0.500 & -- & 0.000 & 0.000 & 0.000 \\
GPT-OSS-20B (think) & 40 & 0.571 & 0.118 & 0.717 & 0.071 & 0.000 & 0.062 & 0.036 & 0.036 & 0.531 & 0.187 & 0.062 & 0.062 & 0.062 \\
GPT-OSS-20B (think) & 50 & 0.440 & 0.222 & 0.562 & 0.182 & 0.000 & 0.149 & 0.098 & 0.096 & 0.481 & -0.053 & 0.182 & 0.182 & 0.000 \\
GPT-OSS-20B (think) & 60 & 0.550 & 0.308 & 0.667 & 0.111 & 0.000 & 0.180 & 0.120 & 0.117 & 0.550 & 0.140 & 0.111 & 0.222 & 0.111 \\
Llama-Nemo-49B (r-on) & 8 & 0.932 & 0.927 & 0.936 & 0.864 & 0.136 & 0.864 & 0.802 & 0.792 & 0.932 & 0.872 & 0.864 & 0.864 & 0.864 \\
Llama-Nemo-49B (r-on) & 10 & 0.909 & 0.900 & 0.917 & 0.818 & 0.227 & 0.818 & 0.751 & 0.740 & 0.909 & 0.832 & 0.818 & 0.818 & 0.818 \\
Llama-Nemo-49B (r-on) & 15 & 0.545 & 0.231 & 0.677 & 0.136 & 0.091 & 0.124 & 0.092 & 0.090 & 0.545 & 0.158 & 0.136 & 0.136 & 0.091 \\
Llama-Nemo-49B (r-on) & 20 & 0.545 & 0.167 & 0.688 & 0.091 & 0.045 & 0.091 & 0.063 & 0.062 & 0.545 & 0.218 & 0.091 & 0.091 & 0.091 \\
Llama-Nemo-49B (r-on) & 30 & 0.545 & 0.444 & 0.615 & 0.182 & 0.000 & 0.264 & 0.209 & 0.205 & 0.545 & 0.098 & 0.182 & 0.364 & 0.182 \\
Llama-Nemo-49B (r-on) & 40 & 0.364 & 0.364 & 0.364 & 0.136 & 0.000 & 0.132 & 0.087 & 0.085 & 0.364 & -0.273 & 0.136 & 0.273 & 0.000 \\
Llama-Nemo-49B (r-on) & 50 & 0.632 & 0.682 & 0.562 & 0.312 & 0.000 & 0.337 & 0.249 & 0.243 & 0.622 & 0.244 & 0.282 & 0.473 & 0.191 \\
Llama-Nemo-49B (r-on) & 60 & 0.409 & 0.381 & 0.435 & 0.000 & 0.000 & 0.165 & 0.112 & 0.110 & 0.409 & -0.183 & 0.000 & 0.364 & 0.000 \\
Nemo-H-47B (r-on) & 8 & 0.668 & 0.594 & 0.720 & 0.338 & 0.002 & 0.346 & 0.293 & 0.288 & 0.668 & 0.362 & 0.338 & 0.356 & 0.337 \\
Nemo-H-47B (r-on) & 10 & 0.641 & 0.638 & 0.645 & 0.301 & 0.001 & 0.302 & 0.251 & 0.246 & 0.641 & 0.283 & 0.301 & 0.321 & 0.286 \\
Nemo-H-47B (r-on) & 15 & 0.582 & 0.482 & 0.650 & 0.193 & 0.000 & 0.202 & 0.160 & 0.157 & 0.583 & 0.179 & 0.193 & 0.227 & 0.171 \\
Nemo-H-47B (r-on) & 20 & 0.556 & 0.406 & 0.646 & 0.144 & 0.000 & 0.153 & 0.121 & 0.118 & 0.556 & 0.131 & 0.144 & 0.173 & 0.130 \\
Nemo-H-47B (r-on) & 30 & 0.537 & 0.346 & 0.641 & 0.116 & 0.000 & 0.126 & 0.097 & 0.095 & 0.537 & 0.090 & 0.116 & 0.152 & 0.093 \\
Nemo-H-47B (r-on) & 40 & 0.505 & 0.067 & 0.663 & 0.021 & 0.000 & 0.022 & 0.017 & 0.016 & 0.505 & 0.029 & 0.021 & 0.028 & 0.017 \\
Nemo-H-47B (r-on) & 50 & 0.503 & 0.030 & 0.666 & 0.009 & 0.000 & 0.009 & 0.007 & 0.007 & 0.503 & 0.027 & 0.009 & 0.010 & 0.008 \\
Nemo-H-47B (r-on) & 60 & 0.500 & -- & 0.667 & 0.000 & 0.000 & 0.000 & 0.000 & 0.000 & 0.500 & -- & 0.000 & 0.000 & 0.000 \\
Qwen3-14B (think) & 8 & 0.853 & 0.829 & 0.871 & 0.706 & 0.473 & 0.708 & 0.638 & 0.629 & 0.853 & 0.737 & 0.706 & 0.711 & 0.706 \\
Qwen3-14B (think) & 10 & 0.755 & 0.680 & 0.802 & 0.517 & 0.351 & 0.514 & 0.445 & 0.436 & 0.757 & 0.585 & 0.515 & 0.518 & 0.515 \\
Qwen3-14B (think) & 15 & 0.568 & 0.478 & 0.632 & 0.311 & 0.096 & 0.297 & 0.233 & 0.227 & 0.570 & 0.150 & 0.311 & 0.385 & 0.177 \\
Qwen3-14B (think) & 20 & 0.516 & 0.604 & 0.379 & 0.225 & 0.000 & 0.212 & 0.157 & 0.153 & 0.516 & 0.037 & 0.225 & 0.290 & 0.093 \\
Qwen3-14B (think) & 30 & 0.501 & 0.660 & 0.066 & 0.030 & 0.000 & 0.034 & 0.020 & 0.020 & 0.501 & 0.007 & 0.030 & 0.035 & 0.028 \\
Qwen3-14B (think) & 40 & 0.500 & 0.657 & 0.079 & 0.035 & 0.000 & 0.041 & 0.023 & 0.023 & 0.500 & 0.000 & 0.035 & 0.043 & 0.023 \\
Qwen3-14B (think) & 50 & 0.539 & 0.693 & 0.072 & 0.039 & 0.000 & 0.038 & 0.021 & 0.021 & 0.503 & 0.015 & 0.039 & 0.039 & 0.026 \\
Qwen3-32B (think) & 8 & 0.824 & 0.788 & 0.850 & 0.647 & 0.482 & 0.648 & 0.576 & 0.566 & 0.823 & 0.687 & 0.645 & 0.655 & 0.645 \\
Qwen3-32B (think) & 10 & 0.692 & 0.564 & 0.762 & 0.389 & 0.308 & 0.393 & 0.325 & 0.318 & 0.692 & 0.474 & 0.389 & 0.399 & 0.384 \\
Qwen3-32B (think) & 15 & 0.580 & 0.370 & 0.685 & 0.225 & 0.096 & 0.224 & 0.175 & 0.171 & 0.581 & 0.220 & 0.223 & 0.245 & 0.173 \\
Qwen3-32B (think) & 20 & 0.548 & 0.475 & 0.604 & 0.284 & 0.046 & 0.272 & 0.210 & 0.205 & 0.550 & 0.104 & 0.278 & 0.389 & 0.141 \\
Qwen3-32B (think) & 30 & 0.497 & 0.586 & 0.361 & 0.214 & 0.000 & 0.199 & 0.144 & 0.141 & 0.497 & -0.006 & 0.214 & 0.283 & 0.059 \\
Qwen3-32B (think) & 40 & 0.505 & 0.667 & 0.041 & 0.021 & 0.000 & 0.021 & 0.012 & 0.012 & 0.505 & 0.043 & 0.021 & 0.021 & 0.016 \\
Qwen3-32B (think) & 50 & 0.583 & 0.722 & 0.167 & 0.091 & 0.000 & 0.091 & 0.052 & 0.052 & 0.545 & 0.227 & 0.091 & 0.091 & 0.091 \\
QwQ-32B (think) & 8 & 0.919 & 0.912 & 0.925 & 0.837 & 0.407 & 0.837 & 0.784 & 0.775 & 0.919 & 0.849 & 0.837 & 0.837 & 0.837 \\
QwQ-32B (think) & 10 & 0.831 & 0.796 & 0.855 & 0.661 & 0.331 & 0.661 & 0.588 & 0.578 & 0.831 & 0.703 & 0.661 & 0.661 & 0.661 \\
QwQ-32B (think) & 15 & 0.598 & 0.339 & 0.711 & 0.206 & 0.115 & 0.204 & 0.155 & 0.151 & 0.598 & 0.316 & 0.206 & 0.206 & 0.196 \\
QwQ-32B (think) & 20 & 0.566 & 0.234 & 0.697 & 0.132 & 0.041 & 0.132 & 0.089 & 0.087 & 0.566 & 0.266 & 0.132 & 0.132 & 0.132 \\
QwQ-32B (think) & 30 & 0.512 & 0.224 & 0.645 & 0.124 & 0.008 & 0.113 & 0.080 & 0.078 & 0.512 & 0.037 & 0.124 & 0.140 & 0.050 \\
QwQ-32B (think) & 40 & 0.559 & 0.554 & 0.565 & 0.309 & 0.000 & 0.298 & 0.233 & 0.228 & 0.560 & 0.121 & 0.309 & 0.436 & 0.155 \\
QwQ-32B (think) & 50 & 0.500 & 0.645 & 0.154 & 0.091 & 0.000 & 0.083 & 0.046 & 0.046 & 0.500 & 0.000 & 0.091 & 0.091 & 0.000 \\
\end{longtable}
\end{landscape}

\section{Architecture-Level Instantiations and Intervention Diagnostics}
\label{sec:supp-model-instantiations}

The main results of this paper are architecture independent. They apply
whenever model behavior is evaluated through binary predictions or scalar
scores. This section illustrates how the same limitations arise in three
familiar model families: additive classifiers, attention-based models, and
graph neural networks (GNNs). These examples do not establish additional
impossibility results. Their purpose is to show why architectural complexity
does not make a higher prediction metric value a reliable indicator of
stronger reasoning.

Across these model families, semantic and non-semantic information can both
affect the final score. If non-semantic information is more predictive on the
benchmark distribution, a model that exploits it can outperform a model that
attempts to capture the intended semantic relationship but does so
imperfectly. The resulting metric ordering is valid as an ordering of
benchmark performance, but it need not be an ordering of reasoning ability.

\subsection{Additive Classifiers over Feature Blocks}
\label{sec:supp-additive-instantiation}

We first consider an additive classifier because the contributions of
semantic and non-semantic information can be represented explicitly. This is
a grouped abstraction of familiar linear, logistic, and additive
models~\cite{joachims1998text,fan2008liblinear,hastie1986generalized,
yuan2006model}.

Following the semantic/non-semantic decomposition in the main paper, let \(z\) denote the
semantic factor that should determine the label, and let \(c\) denote
non-semantic information that may also be available to the model. Let
\[
a_z=\phi_z(z),
\qquad
a_c=\phi_c(c)
\]
be feature vectors derived from these components. Consider a model
\(\model\) with score
\[
s_{\model}(x)
=
w_z^\top a_z
+
w_c^\top a_c
+
b
\]
and binary prediction
\[
\hat y
=
\mathbf{1}[s_{\model}(x)\ge 0].
\]

Define the two score contributions by
\[
S_z(x)=w_z^\top a_z,
\qquad
S_c(x)=w_c^\top a_c.
\]
Then
\[
s_{\model}(x)
=
S_z(x)+S_c(x)+b.
\]

This representation makes two earlier results concrete. First, distinct
semantic and non-semantic contributions can yield the same scalar score and
therefore the same prediction and ranking. This is an instance of the
non-identifiability established in Proposition~1. Second, a larger
non-semantic contribution may improve benchmark predictions or rankings even
when it does not improve semantic reasoning.

For example, consider a semantic classifier and a shortcut classifier with
class-conditional error rates
\[
(\delta_1,\delta_2)
\qquad\text{and}\qquad
(\epsilon_1,\epsilon_2),
\]
respectively. As shown in Proposition~2, the shortcut classifier receives
higher accuracy whenever
\[
\pi\epsilon_1+(1-\pi)\epsilon_2
<
\pi\delta_1+(1-\pi)\delta_2,
\]
where \(\pi=\Pr[Y=1]\). Analogous comparisons hold for balanced accuracy,
G-mean, Youden's \(J\), F1, and other prediction metrics.

The additive model makes the interpretation especially clear. The shortcut
classifier is ranked higher because its non-semantic features produce a more
favorable benchmark error profile. The higher metric value does not imply
that the classifier has captured the semantic relationship that determines
the labels.

A similar conclusion applies to score-based metrics. If the non-semantic
contribution ranks positive and negative benchmark instances more reliably
than the semantic contribution, it can increase AUROC. Thus, observing the
unthresholded score avoids one source of information loss, but it does not
make metric ordering monotone with reasoning ability.

\subsection{Attention-Based Models}
\label{sec:supp-attention-instantiation}

Attention-based models provide a less explicit version of the same
phenomenon. Let an input be
\[
x=(t_1,\ldots,t_n),
\]
where \(t_j\) is the token occurrence at position \(j\). A Transformer maps
the input to hidden representations through repeated attention and
feed-forward transformations~\cite{Vaswani2017Attention}. Abstractly,
\[
H^{(0)}=\operatorname{Embed}(x),
\]
\[
H^{(\ell+1)}
=
\operatorname{Block}_{\ell}(H^{(\ell)}),
\qquad
\ell=0,\ldots,L-1.
\]
After the final layer, the model constructs an instance representation
\[
h_x=\operatorname{Pool}(H^{(L)})
\]
and a scalar score
\[
s_{\model}(x)
=
w_{\mathrm{cls}}^\top h_x+b_{\mathrm{cls}}.
\]

In software-engineering tasks, some input information bears directly on the
intended semantics, while other information may correlate with the labels
without determining them. In vulnerability detection, for example, relevant
semantic information may include relationships among sources, checks,
sanitizers, and sinks. Non-semantic information may include identifier names,
comments, API names, formatting conventions, or dataset-construction
artifacts. In SAT solving, semantic information is distributed across
variables, literals, clauses, and connectives, while generator-specific
formatting or structural regularities may provide non-semantic cues.

Attention-based models can use either kind of information. A model that
captures the relevant semantic relationships may still make errors because
those relationships are difficult to learn. Another model may exploit easier
non-semantic correlations and obtain a better benchmark error profile or
ranking. In that case,
\[
M(B,\model_c)>M(B,\model_z)
\]
may hold even though \(\model_z\) relies more strongly on the intended
semantic relationship.

This is the practically important implication of the attention
instantiation. A higher accuracy, F1, ADR, or AUROC does not establish that
the higher-scoring attention model reasons more effectively. It establishes
only that the model performs better according to the observable criterion of
the metric on the evaluated distribution.

Attention weights, gradients, and attribution methods such as Integrated
Gradients, SHAP, and DeepLIFT can provide additional diagnostic
information~\cite{Sundararajan2017IG,LundbergLee2017SHAP,
Shrikumar2017DeepLIFT}. However, these analyses should not be treated as
direct measures of reasoning. Their results depend on the explanation method,
the selected representation, and how interactions among tokens are
allocated. They can motivate further tests, but they do not by themselves
justify interpreting a higher prediction metric as stronger reasoning.

\subsection{Graph Neural Networks and Message Passing}
\label{sec:supp-gnn-instantiation}

Graph neural networks instantiate the same limitation in models that operate
over structured program representations. Let
\[
G=(V,E,X)
\]
be a program graph, where \(V\) is the node set, \(E\) is the edge set, and
\(X_v\) is the initial feature vector of node \(v\). A generic
message-passing layer computes
\[
m_v^{(k)}
=
\Agg_{u\in N(v)}
\Msg_k(h_v^{(k)},h_u^{(k)},e_{uv}),
\]
followed by
\[
h_v^{(k+1)}
=
\Upd_k(h_v^{(k)},m_v^{(k)}).
\]
After \(K\) layers, the model produces a graph representation
\[
h_G
=
\Readout_{v\in V}h_v^{(K)}
\]
and score
\[
s_{\model}(G)=w^\top h_G+b.
\]

Program graphs may encode AST, control-flow, data-flow, call, or dependence
relationships. The semantic condition underlying a vulnerability may depend
on a long-range relation, such as an untrusted value reaching a sink without
sanitization, a missing bounds check before a memory access, or a particular
alias being dereferenced under an unsafe path. These relationships can be
difficult to represent and learn.

At the same time, graph inputs may contain easier non-semantic correlations,
including risky API names, identifier features, project-specific graph
patterns, or local motifs introduced by data extraction. A GNN that exploits
these correlations may obtain a higher metric value than a detector that
attempts to capture the relevant source-to-sink or control-flow relationship
but does so imperfectly.

The resulting ranking is again non-monotone with respect to reasoning
ability. A higher metric value can indicate that local or project-specific
features are more predictive on the benchmark. It does not necessarily
indicate that the model has captured the vulnerability condition represented
by the graph.

Graph explanation methods can help diagnose this possibility. For example,
GNNExplainer identifies a subgraph and feature subset that approximately
preserve a model's prediction~\cite{Ying2019GNNExplainer}. An explanation
centered on a vulnerability-relevant slice may provide encouraging evidence,
whereas one centered on unrelated names or local graph motifs may reveal a
potential shortcut. Nevertheless, an explanation describes what influenced
the model's output under the assumptions of the explanation method. It does
not establish that the model's metric advantage reflects stronger
vulnerability reasoning.

\subsection{Intervention-Based Comparison}
\label{sec:supp-instantiation-interventions}

The model-level examples suggest a common way to strengthen reasoning-oriented
evaluation. Rather than interpreting an aggregate metric value in isolation,
an evaluator can examine how model outputs respond to controlled semantic and
non-semantic changes.

Let \(x_{\Delta z}\) be an instance in which the semantic condition is changed
while non-semantic information is held as fixed as possible. Let
\(x_{\Delta c}\) preserve the semantic condition while changing
non-semantic information. Define the score changes
\[
\Delta_z(x)
=
\left|
s_{\model}(x)-s_{\model}(x_{\Delta z})
\right|
\]
and
\[
\Delta_c(x)
=
\left|
s_{\model}(x)-s_{\model}(x_{\Delta c})
\right|.
\]

A model whose outputs track the intended task semantics should generally
respond to semantic changes that alter the correct answer and remain
comparatively stable under semantics-preserving non-semantic changes. For
binary predictions, analogous quantities can be defined through prediction
flips:
\[
F_z(x)
=
\mathbf{1}
[
\model(x)\neq\model(x_{\Delta z})
],
\]
\[
F_c(x)
=
\mathbf{1}
[
\model(x)\neq\model(x_{\Delta c})
].
\]

Matched SAT/UNSAT formulas and vulnerable/fixed program variants approximate
semantic interventions. Identifier renaming, formatting changes, irrelevant
code transformations, and generator perturbations approximate
semantics-preserving non-semantic interventions. Pair, group, and cluster
metrics summarize success over such controlled instances.

These interventions provide a more informative evaluation design, but the
limitations of controlled metrics established in the main paper still apply.
A higher controlled score indicates stronger performance under the selected
interventions and aggregation rule. It does not automatically imply stronger
reasoning, because residual non-semantic correlations, intervention quality,
dependence among instances, and threshold selection can still affect the
ordering.

\subsection{Summary}
\label{sec:supp-instantiation-summary}

The three instantiations differ in how information is represented and
combined, but they support the same interpretation. Additive classifiers may
combine explicit feature-block contributions. Attention-based models mix
information across token representations. GNNs aggregate node and edge
information through message passing. These architectural differences can
help explain why a model succeeds, but they do not change what a
prediction-output metric establishes.

Across all three settings, a higher metric value establishes better
performance under the metric's observable benchmark criterion. It does not,
without controlled interventions or additional task-specific evidence,
establish stronger reasoning. Architecture-specific explanations and
attributions can provide useful diagnostics, while witnesses and controlled
semantic and non-semantic interventions provide stronger behavioral
evidence. None of these observations makes an ordinary prediction metric a
standalone measure of reasoning ability.

\section{Reproduction Notes and Remaining Requirements}
\label{sec:supp-reproduction}

The accompanying revision package includes CSV files transcribed from the
source LaTeX tables and CSV outputs for the metric-level, ordering,
pair-rematching, and confidence-bound analyses. The strict-ordering
calculation should be rerun on full-precision values before submission.
Per-instance predictions should then be used to verify every source-table
cell.

The source matched-separability protocol describes a matched-pair cluster
bootstrap for dependence within seed families. The supplied summary files do
not contain per-instance outcomes or bootstrap replicates, so this revision
does not report invented bootstrap values. A final artifact should resample
matched pairs within groups, recompute the group-averaged
\(\hat{s}\), and report a dependence-robust lower quantile beside the
Clopper--Pearson and Wilson reference bounds.

Witness validation should also be strengthened on the negative side.
SAT experiments should elicit and check a normalized UNSAT proof when
feasible; graph-coloring experiments should use a checkable non-colorability
certificate or a trusted solver trace. Until then, \(W\) remains a
positive-side constructive reference rather than a complete semantic oracle.

\bibliographystyle{ACM-Reference-Format}
\bibliography{paper}

\end{document}
